\SourceLeaf{35}
\fancyhead[L]{\small Allgemeine Grundsätze der Theorie der Gestalt von Flüssigkeiten}
\renewcommand{\Article}[1]{\par\medskip\Needspace{4\baselineskip}\pdfbookmark[1]{Art. #1}{fluidorum-art-#1}\centerline{\textbf{#1.}}\nobreak\smallskip}
\pdfbookmark[0]{Allgemeine Grundsätze der Theorie der Gestalt von Flüssigkeiten im Gleichgewicht}{fluidorum-body}
\par\medskip
\centerline{ALLGEMEINE GRUNDSÄTZE}
\centerline{DER THEORIE DER GESTALT VON FLÜSSIGKEITEN}
\centerline{IM GLEICHGEWICHT.}
\medskip
\expandafter\def\csname HistoricalNoteText4\endcsname{Es steht fest, dass sich die größte Anziehung, die eine gegebene homogene Masse nach jenem Gesetz auf einen gegebenen Punkt ausüben kann, zu der Anziehung, die dieselbe zu einer Kugel geformte Masse auf einen Punkt an ihrer Oberfläche ausübt, wie $3$ zu $\BareRoot[3]{25}$ verhält. Die letztere Anziehung lässt sich leicht mit der Schwerkraft vergleichen.}
\expandafter\def\csname HistoricalNoteText5\endcsname{So sind die meisten Einwände zu beurteilen, die in der Zeitschrift von Pavia (Giornale di fisica etc. T. 9) vorgebracht wurden. Der angesehene \textsc{Petit} hat ihnen in den Annales de chimie et de physique T. 4 treffend geantwortet.}
Der scharfsinnige \textsc{Clairaut} hat als Erster mit Scharfblick und Genauigkeit die Kräfte aufgezählt, die den Anstieg oder das Absinken von Flüssigkeiten in Kapillarröhren bestimmen. Da er das Gesetz dieser Kräfte aber völlig unberührt ließ, konnte diese Aufzählung nichts zur mathematischen Erklärung der Erscheinungen beitragen. Die gewöhnliche Anziehung, die dem Quadrat der Entfernung umgekehrt proportional ist und alle Himmelsbewegungen so erfolgreich erklärt, ist für die Erklärung der Kapillarerscheinungen sowie der Adhäsions- und Kohäsionserscheinungen unbrauchbar. Eine richtig durchgeführte Rechnung zeigt leicht, dass nach diesem Gesetz die Anziehung jedes Körpers, mit dem sich Experimente anstellen lassen, dessen Größe also gegenüber der ganzen Erde als verschwindend gelten kann, an einem beliebigen Punkt gegenüber der Schwerkraft verschwindet, selbst wenn der Punkt den Körper berührt\HistoricalNote{4}. Daraus ist mit Recht zu schließen, dass dieses Anziehungsgesetz bei kleinsten Entfernungen mit der Natur nicht mehr übereinstimmt und eine gewisse Abänderung verlangt. Gleichbedeutend ist, dass die Teilchen der Körper neben jener Anziehungskraft eine weitere Kraft ausüben, die nur bei kleinsten Entfernungen merklich wird. Alle Erscheinungen deuten übereinstimmend darauf hin, dass dieser zweite Anteil der Anziehungskraft (die \textit{Molekular\SourceLeaf{36}anziehung}) selbst bei den kleinsten messbaren Entfernungen unmerklich ist, während er bei unmerklich kleinen Entfernungen den ersten, dem Quadrat der Entfernung umgekehrt proportionalen Anteil weit übertreffen kann.

Der bedeutende \textsc{Laplace} ging von dieser einzigen Annahme über die Natur der Molekularkräfte aus und ließ das Gesetz ihrer Abnahme bei zunehmenden Entfernungen im Übrigen völlig unbestimmt. Als Erster berechnete er genau, wie diese Kräfte die Gestalt der Oberfläche von Flüssigkeiten beeinflussen. Nachdem er eine allgemeine Gleichung für die Gleichgewichtsgestalt aufgestellt hatte, versuchte er daraus sowohl die Kapillarerscheinungen im eigentlichen Sinn als auch viele verwandte Erscheinungen zu erklären. Die erstaunliche Übereinstimmung mit sorgfältigen Experimenten hat diese Untersuchungen durchweg bestätigt. Sie zählen zu den schönsten Fortschritten der Naturwissenschaft, die wir jenem großen Geometer verdanken. Die Einwände einiger Autoren gegen diese Untersuchungen sind größtenteils von geringer oder gar keiner Bedeutung\HistoricalNote{5}.

In den Rechnungen des bedeutenden \textsc{Laplace} finden sich allerdings einige Stellen, die mit einer strengen Beweisführung nicht ganz vereinbar sind. In der früheren Abhandlung \textit{théorie de l’action capillaire} bezeichnet $\varphi f$ die Stärke der Anziehung in der Entfernung $f$. Das Integral $\int\varphi f.\mathrm{d}f$ wird von $f=x$ bis $f=\infty$ erstreckt und $=\Pi x$ gesetzt. Dann wird das Integral $\int\Pi f.f\mathrm{d}f$ von $f=x$ bis $f=\infty$ genommen und $=\Psi x$ gesetzt. Schließlich werden die Werte der Integrale $2\pi\int\Psi f.\mathrm{d}f$, $2\pi\int\Psi f.f\mathrm{d}f$ von $f=0$ bis $f=\infty$ der Reihe nach $=K$ und $=H$ gesetzt, wobei $\pi$ den halben Umfang eines Kreises mit Radius $=1$ bezeichnet. Die Beschaffenheit der Funktion $\varphi f$ bleibt völlig offen, sofern sie für alle merklichen Werte von $f$ unmerklich ist. Aus dieser Annahme \textit{allein} würde aber keineswegs folgen, dass auch $\Pi f,\Psi f$ für merkliche Werte von $f$ notwendig unmerklich werden. Ebenso wenig darf man daraus schließen, dass sich die Integrale $2\pi\int\Psi f.\mathrm{d}f$, $2\pi\int\Psi f.f\mathrm{d}f$, erstreckt von $f=0$ bis zu einem merklichen \textit{endlichen} Wert von $f$, nur unmerklich von $K,H$ unterscheiden, wie es in jener Abhandlung heißt. Denn es ließen sich unendlich viele Formen der Funktion $\varphi f$ denken, die der Grundannahme genügen und für die diese Schlüsse falsch wären. Nimmt man sogar an, dass $\varphi f$ die gesamte Anziehung ausdrückt, so enthält die Funktion tatsächlich auch einen Anteil der Form $\frac{a}{ff}$, von dem die gewöhnliche Anziehung abhängt. Selbst wenn aber dieser Term als unmerklich gelten muss, solange die Ab\SourceLeaf{37}messungen der anziehenden Körper, wie sie in Experimenten vorkommen können, gegenüber der ganzen Erde unmerklich sind, würde bereits die zweite Integration bei einer Erstreckung ins Unendliche einen unendlichen Term in die Funktion $\Psi f$ einführen.

Wenn sich in diesen und ähnlichen Punkten eine leichte Nachlässigkeit anzudeuten scheint, so betrifft sie gewiss eher die Art der Darstellung als die Sache selbst. Aus der zweiten Abhandlung \textit{Supplément à la théorie de l’action capillaire} geht nämlich hervor, dass der bedeutende \textsc{Laplace} unter $\varphi f$ stillschweigend nur den Anteil verstand, der zur gewöhnlichen Anziehung hinzukommt, und nicht die gesamte Anziehung. Die gewöhnliche Anziehung kann, wie leicht einzusehen ist, in unseren Experimenten keine merkliche Veränderung bewirken. Er deutet sogar an, dass er die Funktion $\varphi f$ nach Art der Exponentialfunktion $e^{-if}$ betrachtet, wobei $i$ einen sehr großen Wert oder vielmehr $\frac{1}{i}$ eine sehr kleine Länge bezeichnet. Es ist aber nicht einmal nötig, die Allgemeinheit so stark einzuschränken. Wer auf die Sache und weniger auf die Worte achtet, sieht leicht, dass es genügt, diese Integrationen nur bis zu einer beliebigen merklichen Entfernung statt bis ins Unendliche zu erstrecken, oder, wenn man so will, bis zu einer endlichen Entfernung, die größer ist als die in den Experimenten vorkommenden Abmessungen.

Diese Theorie leidet aber an einem anderen, weit schwerwiegenderen Mangel, den ihre Kritiker, soweit wir wissen, nicht einmal bemerkt haben. Die Theorie besteht aus zwei Teilen. Der eine stellt für die freie Flüssigkeitsoberfläche eine allgemeine Gleichung zwischen den partiellen Differenzialen der Koordinaten auf. Die Gleichung hängt von der molekularen Anziehungskraft ab, die die Flüssigkeitsteilchen gegenseitig aufeinander ausüben. Dieser Teil der Theorie ist so vollständig ausgearbeitet, dass nichts Wesentliches zu wünschen übrig bleibt. Eine solche Gleichung zwischen partiellen Differenzialen, deren Integration, wenn die Analysis sie leisten könnte, willkürliche Funktionen einführen würde, genügt aber nicht, um die Gestalt der Oberfläche \textit{vollständig} zu bestimmen. Dazu muss eine \textit{neue} Bedingung für die Beschaffenheit der Gestalt an den Grenzen hinzukommen und eine solche Bedingung stellt der zweite Teil der Theorie auf. Danach muss der Winkel zwischen der Tangentialebene an die freie Flüssigkeitsoberfläche am Rand des Gefäßes, genauer an der Grenze der merklichen Anziehungskraft der Gefäßwand, und der Tangentialebene an die Gefäßwand an derselben Stelle \textit{konstant} sein. Bestimmt wird dieser Winkel durch die Beziehung zwischen den Stärken der Molekularkräfte des Gefäßes und der Flüssigkeit, sofern die Gestalt des Gefäßes am Rand der freien Flüssigkeitsoberfläche stetig ist. Doch der bedeutende \textsc{Laplace} hat nicht einmal versucht, diesen zentralen Satz der ganzen Theorie durch Rechnung zu beweisen. Die hierzu auf S. 5 der früheren Abhandlung vorge\SourceLeaf{38}brachten Überlegungen liefern nur eine vage Argumentation und setzen bereits voraus, was zu beweisen war. Die auf S. 44 und den folgenden Seiten angestellten Rechnungen bleiben ohne Ergebnis. In der zweiten Abhandlung wird zwar der Anstieg der Flüssigkeit in Kapillarröhren nach einer anderen Methode behandelt. Vergleicht man ihr Ergebnis mit dem der früheren Methode, so erhält man eine Formel für jenen Winkel zwischen den Tangentialebenen, und zwar eine richtige. Es ist aber zu beachten, dass hier streng genommen die Konstanz des Winkels bereits \textit{vorausgesetzt} wird. Außerdem ist die an sich wenig befriedigende Methode auf den sehr speziellen Fall eines prismatischen Gefäßes mit senkrechten Wänden beschränkt. Angesichts dessen muss man einräumen, dass die vom bedeutenden \textsc{Laplace} aufgestellte Theorie in wesentlichen Punkten noch Mängel aufweist und unvollständig ist.

Wir werden daher die Theorie der Gleichgewichtsgestalt von Flüssigkeiten unter der Wirkung der Schwerkraft sowie der Molekularkräfte der Flüssigkeit selbst und des Gefäßes von Grund auf neu entwickeln. Dabei werden wir einer ganz anderen, aus den ersten Prinzipien der Dynamik gewonnenen Methode folgen und von Anfang an die größtmögliche Allgemeinheit wählen. Diese Untersuchung wird zu einem bemerkenswerten neuen Satz führen, der die vollständige Theorie in einer einzigen, überaus einfachen Formel zusammenfasst. Aus ihm ergeben sich beide Teile der Theorie des bedeutenden \textsc{Laplace} von selbst.
\par\medskip\centerline{\rule{.3\textwidth}{.3pt}}

